\documentclass[12pt,a4paper]{article}

\usepackage[utf8]{inputenc}
\usepackage[margin=1in]{geometry}
\usepackage{xcolor}
\usepackage{tcolorbox}
\usepackage{listings}
\usepackage{hyperref}

\definecolor{openglblue}{RGB}{85,134,164}
\definecolor{codegray}{RGB}{240,240,240}
\definecolor{strgray}{RGB}{96,96,96}
\definecolor{kwblue}{RGB}{34,113,179}

\tcbuselibrary{skins, breakable}

\newtcolorbox{functionbox}[2][]{
    enhanced,
    breakable,
    colback=codegray,
    colframe=openglblue,
    title=\textbf{#2},
    fonttitle=\sffamily,
    attach boxed title to top left={yshift=-2mm, xshift=5mm},
    boxed title style={colback=openglblue},
    #1
}

\lstset{
    backgroundcolor=\color{codegray},
    commentstyle=\color{strgray}\itshape,
    keywordstyle=\color{kwblue}\bfseries,
    numberstyle=\tiny\color{strgray},
    basicstyle=\ttfamily\footnotesize,
    breakatwhitespace=false,         
    breaklines=true,                 
    captionpos=b,                    
    keepspaces=true,                 
    showspaces=false,                
    showstringspaces=false,
    showtabs=false,                  
    tabsize=4,
    language=C++
}

\title{\textbf{Building a Mini Game Engine with OpenGL}\\ \large Complete Requirements \& Implementation Guide}
\author{Beginner's Guide}
\date{}

\begin{document}

\maketitle

\tableofcontents
\newpage

\section*{The Big Picture}
Before diving in: OpenGL is a \textbf{state machine API} that talks to your GPU. You write C++ code that sends data and instructions to the GPU driver, which compiles shaders and executes them on thousands of tiny parallel cores. The CPU sets everything up; the GPU does the actual pixel math.

\vspace{1em}
This document expands on all 11 requirements, explaining not just what they are, but \textbf{how to implement them} in C++ code using OpenGL.

\section{Requirement 1: Shader Program}

\subsection*{What it is}
A shader is a small program that runs \textit{on the GPU}. The \textbf{vertex shader} runs once per vertex and decides \textit{where} the point lands on screen. The \textbf{fragment shader} runs once per pixel and decides \textit{what color} it gets.

\subsection*{Implementation Details}
To implement this, you read GLSL code (a C-like language) from a file, compile it, and link it into an active shader program.

\begin{lstlisting}[language=C++]
// 1. Upload source code to GPU
GLuint vertexShader = glCreateShader(GL_VERTEX_SHADER);
glShaderSource(vertexShader, 1, &vertexSourceString, NULL);

// 2. Compile on the GPU
glCompileShader(vertexShader);

// 3. Link Vertex + Fragment into a Program
GLuint shaderProgram = glCreateProgram();
glAttachShader(shaderProgram, vertexShader);
glAttachShader(shaderProgram, fragmentShader);
glLinkProgram(shaderProgram);

// 4. Activate it before drawing
glUseProgram(shaderProgram);
\end{lstlisting}


\section{Requirement 2: Mesh}

\subsection*{What it is}
A mesh is the 3D geometry of an object—a collection of vertices (points in 3D space) and faces (triangles connecting those points).

\subsection*{Implementation Details}
You need to configure three GPU objects to successfully draw a mesh.
\begin{itemize}
    \item \textbf{VBO (Vertex Buffer Object):} A GPU array holding vertex data.
    \item \textbf{EBO (Element Buffer Object):} A GPU array holding indices (which triangles connect which vertices).
    \item \textbf{VAO (Vertex Array Object):} A config object that tells the shaders how to interpret the VBO bytes.
\end{itemize}

\begin{lstlisting}[language=C++]
// Create the objects
GLuint VAO, VBO, EBO;
glGenVertexArrays(1, &VAO);
glGenBuffers(1, &VBO);
glGenBuffers(1, &EBO);

// Bind VAO first, then upload data to VBO and EBO
glBindVertexArray(VAO);

glBindBuffer(GL_ARRAY_BUFFER, VBO);
glBufferData(GL_ARRAY_BUFFER, sizeof(vertices), vertices, GL_STATIC_DRAW);

glBindBuffer(GL_ELEMENT_ARRAY_BUFFER, EBO);
glBufferData(GL_ELEMENT_ARRAY_BUFFER, sizeof(indices), indices, GL_STATIC_DRAW);

// Tell the VAO how to read positions (e.g., 3 floats at location 0)
glVertexAttribPointer(0, 3, GL_FLOAT, GL_FALSE, 8 * sizeof(float), (void*)0);
glEnableVertexAttribArray(0);

// Later, to draw:
glBindVertexArray(VAO);
glDrawElements(GL_TRIANGLES, 6, GL_UNSIGNED_INT, 0);
\end{lstlisting}


\section{Requirement 3: Transform}

\subsection*{What it is}
A Transform encodes an object's position, rotation, and scale into a single $4\times4$ matrix.

\subsection*{Implementation Details}
Using a math library like \texttt{GLM}, you calculate the $4\times4$ matrix locally on the CPU, and then send it to the GPU via a \textbf{uniform}.

\begin{lstlisting}[language=C++]
#include <glm/glm.hpp>
#include <glm/gtc/matrix_transform.hpp>

// Create a transform matrix (Translation * Rotation * Scale)
glm::mat4 model = glm::mat4(1.0f);
model = glm::translate(model, glm::vec3(5.0f, 1.0f, -3.0f)); // Move
model = glm::rotate(model, glm::radians(45.0f), glm::vec3(0,1,0)); // Rotate Y
model = glm::scale(model, glm::vec3(0.5f)); // Scale to 50%

// Upload Matrix to the Vertex Shader Uniform
GLint transformLoc = glGetUniformLocation(shaderProgram, "transform");
glUniformMatrix4fv(transformLoc, 1, GL_FALSE, &model[0][0]);
\end{lstlisting}

Inside the vertex shader, you then apply this: \lstinline!gl_Position = transform * vec4(position, 1.0);!


\section{Requirement 4: Pipeline State}

\subsection*{What it is}
OpenGL needs you to flip specific state switches before drawing—such as depth testing (checking what objects are in front of others) and blending (glass transparency). 

\subsection*{Implementation Details}
You'll structure a class that encapsulates states so you can quickly toggle them before drawing.
\begin{lstlisting}[language=C++]
void PipelineState::setup() {
    // Depth Testing (ensure geometry doesn't draw through walls)
    if (depthTestEnabled) {
        glEnable(GL_DEPTH_TEST);
        glDepthFunc(GL_LESS);
    } else {
        glDisable(GL_DEPTH_TEST);
    }

    // Blending (for glass/transparent objects)
    if (blendingEnabled) {
        glEnable(GL_BLEND);
        glBlendFunc(GL_SRC_ALPHA, GL_ONE_MINUS_SRC_ALPHA);
    } else {
        glDisable(GL_BLEND);
    }

    // Face Culling (skip drawing inside of models)
    if (cullFaceEnabled) glEnable(GL_CULL_FACE);
    else glDisable(GL_CULL_FACE);
}
\end{lstlisting}


\section{Requirement 5: Texture}

\subsection*{What it is}
A texture is an image (PNG/JPG) uploaded to the GPU to give models detail (like bricks or wood). 

\subsection*{Implementation Details}
1. Load image via CPU (\texttt{stb\_image}). 2. Upload to GPU to a Texture ID.
\begin{lstlisting}[language=C++]
// 1. Generate texture ID
GLuint texture;
glGenTextures(1, &texture);
glBindTexture(GL_TEXTURE_2D, texture);

// 2. Load pixels using STB
int width, height, channels;
unsigned char *data = stbi_load("bricks.png", &width, &height, &channels, 0);

// 3. Upload to GPU
if (data) {
    glTexImage2D(GL_TEXTURE_2D, 0, GL_RGB, width, height, 0, 
                 GL_RGB, GL_UNSIGNED_BYTE, data);
    stbi_image_free(data);
}
\end{lstlisting}


\section{Requirement 6: Sampler}

\subsection*{What it is}
A sampler defines how a texture is filtered. When the image is stretched, does it get blurry (Linear) or blocky (Nearest/Pixelated)?

\subsection*{Implementation Details}
You set Sampler parameters directly on the \texttt{GL\_TEXTURE\_2D} object or via dedicated Sampler Objects.
\begin{lstlisting}[language=C++]
glBindTexture(GL_TEXTURE_2D, texture);

// How to sample when stretching (MAG_FILTER) and shrinking (MIN_FILTER)
glTexParameteri(GL_TEXTURE_2D, GL_TEXTURE_MIN_FILTER, GL_LINEAR_MIPMAP_LINEAR);
glTexParameteri(GL_TEXTURE_2D, GL_TEXTURE_MAG_FILTER, GL_LINEAR);

// How to handle UVs outside 0.0 - 1.0 (Wrap/Repeat)
glTexParameteri(GL_TEXTURE_2D, GL_TEXTURE_WRAP_S, GL_REPEAT);
glTexParameteri(GL_TEXTURE_2D, GL_TEXTURE_WRAP_T, GL_REPEAT);

// Generate Mipmaps (smaller versions of texture for objects far away)
glGenerateMipmap(GL_TEXTURE_2D);
\end{lstlisting}


\section{Requirement 7: Material}

\subsection*{What it is}
A Material bundles a Shader, Pipeline State, and Textures into one object describing "how this should look."

\subsection*{Implementation Details}
You write a generic \texttt{Material} class with a \texttt{setup()} function called right before a draw call.
\begin{lstlisting}[language=C++]
class TexturedMaterial : public Material {
public:
    ShaderProgram* shader;
    PipelineState pipeline;
    Texture* diffuseTex;
    
    void setup() override {
        // 1. Activate Material's Shader
        shader->use();
        
        // 2. Setup States (Depth, Blend, etc.)
        pipeline.setup();
        
        // 3. Bind Textures and Uniforms
        glActiveTexture(GL_TEXTURE0);
        glBindTexture(GL_TEXTURE_2D, diffuseTex->id);
        shader->setInt("sampler", 0); // Tell shader texture is in slot 0
    }
};
\end{lstlisting}


\section{Requirement 8: Entities and Components (ECS)}

\subsection*{What it is}
An architectural pattern organizing your game world by composing objects from reusable data pieces, avoiding huge rigid class structures.

\subsection*{Implementation Details}
An \texttt{Entity} is just an ID. It holds components like a \texttt{TransformComponent} and a \texttt{MeshRendererComponent}.
\begin{lstlisting}[language=C++]
struct TransformComponent {
    glm::vec3 position;
    glm::vec3 rotation;
    glm::vec3 scale;
};

struct MeshRendererComponent {
    Mesh* mesh;
    Material* material;
};

// In your world structure:
class World {
    std::vector<Entity*> entities;
    // Systems loop over entities and draw them if they have 
    // a mesh renderer and transform component.
};
\end{lstlisting}


\section{Requirement 9: Forward Renderer}

\subsection*{What it is}
The system that loops through all entities and draws them in the correct order. Transparent objects MUST be drawn \textbf{after} opaque objects and sorted from back to front.

\subsection*{Implementation Details}
\begin{lstlisting}[language=C++]
void renderScene(World& world) {
    // 1. Reset Screen and Depth
    glClear(GL_COLOR_BUFFER_BIT | GL_DEPTH_BUFFER_BIT);

    // 2. Draw Opaque (Solid) Objects
    for (Entity* e : world.opaqueEntities) {
        e->meshRenderer.material->setup();
        uploadTransform(e->transform);
        e->meshRenderer.mesh->draw(); // Calls glDrawElements
    }

    // 3. Sort Transparent Objects by distance to Camera 
    sortByDistanceToCamera(world.transparentEntities);

    // 4. Draw Transparent Objects (with blending enabled in material)
    for (Entity* e : world.transparentEntities) {
        e->meshRenderer.material->setup();
        uploadTransform(e->transform);
        e->meshRenderer.mesh->draw();
    }
}
\end{lstlisting}


\section{Requirement 10: Sky Rendering}

\subsection*{What it is}
Drawing a massive sky sphere texture that never gets closer or further away.

\subsection*{Implementation Details}
1. \textbf{In C++}: Strip translation (movement) from the Camera's View Matrix.
\begin{lstlisting}[language=C++]
// Convert View Mat4 -> Mat3 (Removes translation vector) -> back to Mat4
glm::mat4 skyView = glm::mat4(glm::mat3(cameraViewMatrix));
glUniformMatrix4fv(viewLoc, 1, GL_FALSE, &skyView[0][0]);
\end{lstlisting}

2. \textbf{In Vertex Shader}: Force the Z value to exactly 1.0 (maximum depth).
\begin{lstlisting}[language=C++]
// GLSL Vertex Shader trick:
vec4 pos = projection * view * vec4(position, 1.0);
// Set Z (depth) equal to W so perspective division (Z/W) always equals 1.0
gl_Position = pos.xyww; 
\end{lstlisting}


\section{Requirement 11: Post-Processing}

\subsection*{What it is}
Drawing to a custom hidden image (Framebuffer), then applying full-screen 2D effects (color grading, vignette) to it.

\subsection*{Implementation Details}
You use a Framebuffer Object (FBO).
\begin{lstlisting}[language=C++]
// SETUP: Create Framebuffer
GLuint FBO, textureColorBuffer;
glGenFramebuffers(1, &FBO);
glBindFramebuffer(GL_FRAMEBUFFER, FBO);

// Attach a blank texture to FBO to act as our "Screen"
glGenTextures(1, &textureColorBuffer);
glBindTexture(GL_TEXTURE_2D, textureColorBuffer);
// ...allocate texture space...
glFramebufferTexture2D(GL_FRAMEBUFFER, GL_COLOR_ATTACHMENT0, 
                       GL_TEXTURE_2D, textureColorBuffer, 0);

// RENDER LOOP:
// 1. Render to Custom Buffer
glBindFramebuffer(GL_FRAMEBUFFER, FBO);
glClear(GL_COLOR_BUFFER_BIT | GL_DEPTH_BUFFER_BIT);
renderScene(world); 

// 2. Render Custom Buffer to Native Screen with Post-Process shader
glBindFramebuffer(GL_FRAMEBUFFER, 0); 
glClear(GL_COLOR_BUFFER_BIT);
postProcessShader.use();
glBindTexture(GL_TEXTURE_2D, textureColorBuffer);
drawFullScreenQuad();
\end{lstlisting}

\end{document}
